\section*{الفصل \ref{ch:b2:corrosion} --- التآكل والحماية}

\begin{solution}{exo:b2:corrosion:1}
$j = \qty{0.020}{A/m^2}$: $de/dt = 0.020 \times 65.38/(2 \times 96\,485 \times 7.13
\times 10^6) = \qty{9.5e-13}{m/s}$، أي \qty{0.030}{mm} في السنة.
\end{solution}

\begin{solution}{exo:b2:corrosion:2}
براغي الفولاذ على النحاس: الحديد، الأقل نبلًا، هو المصعد، والبراغي الصغيرة
على مهبط نحاسي كبير تتآكل بسرعة. ومسامير الزنك في الفولاذ: يتآكل الزنك
ويحمي الفولاذ من حوله.
\end{solution}

\begin{solution}{exo:b2:corrosion:3}
في الفجوة الضيقة بين الصفيحتين وتحت رأس البرغي، حيث يتجدد الماء
ببطء ويُستنفد أكسجينه سريعًا: فبفعل \omterm{def:b2:corrosion:differential}{التهوية التفاضلية}
تصير هذه المناطق المحصورة مصاعد، والسطوح المكشوفة مهابط.
\end{solution}

\begin{solution}{exo:b2:corrosion:4}
العلبة المقصدرة: فالقصدير، الأنبل من الحديد، يجعل الفولاذ المكشوف مصعدًا لمهبط
كبير. أما في الدلو المجلفن فيتآكل الزنك بدل
الفولاذ.
\end{solution}

\begin{solution}{exo:b2:corrosion:5}
التحريك يُرقّق \omterm{def:b2:current-potential-curves:limiting-current}{طبقة الانتشار}: فيرتفع استواء الأكسجين، ويلاقيه المنحنى المصعدي
للحديد أعلى، فيزداد \omterm{def:b2:corrosion:uniform}{جهد التآكل} \omterm{def:b2:corrosion:uniform}{وتيار التآكل}
كلاهما. ودون أكسجين، لا يبقى من تفاعل مهبطي في الماء
المتعادل إلا الاختزال البطيء للماء: فينخفض \omterm{def:b2:corrosion:uniform}{جهد التآكل} ويصير
\omterm{def:b2:corrosion:uniform}{تيار التآكل} صغيرًا جدًا.
\end{solution}

\begin{solution}{exo:b2:corrosion:6}
$t = 0.90 \times 5000 \times 2 \times 96\,485/(65.38 \times 0.10) = \qty{1.33e8}{s}$،
أي نحو 4.2~سنة.
\end{solution}

\begin{solution}{exo:b2:corrosion:7}
الشحنة لكل كيلوغرام، $nF/M$: الزنك \qty{820}{A.h/kg}، والمغنيسيوم \qty{2.21e3}{A.h/kg}،
والألومنيوم \qty{2.98e3}{A.h/kg}. والجهود المعيارية: $\ce{Zn^2+}/\ce{Zn}$ $-0.76$،
$\ce{Mg^2+}/\ce{Mg}$ $-2.36$، $\ce{Al^3+}/\ce{Al}$ \qty{-1.68}{V}. وفي تربة جافة مقاوِمة
يجب دفع التيار عبر مقاومة كبيرة: فالمغنيسيوم، الواقع أدنى بكثير
من الحديد، يقدم أكبر جهد دافع.
\end{solution}

\begin{solution}{exo:b2:corrosion:8}
$U = 2.0 \times 4.0 + 1.5 = \qty{9.5}{V}$؛ الاستطاعة \qty{19}{W}؛ وفي السنة $19 \times
8766 = \qty{1.7e5}{W.h}$، أي \qty{1.7e2}{kW.h}.
\end{solution}

\begin{solution}{exo:b2:corrosion:9}
عند خط الماء يبلل الفلزَ ماءٌ غني بالأكسجين، فيكون مهبطًا؛ وتحته
بقليل يكون الماء أقل تهوية فتصير المنطقة مصعد الخلية
المتكوّنة مع خط الماء: فيتآكل الفولاذ هناك، في شريط.
\end{solution}

\begin{solution}{exo:b2:corrosion:10}
في الماء المهوّى يقع جهد تآكل الفولاذ المقاوم للصدأ على استواء
تخميله: تيار ضئيل، والغشاء يرمم نفسه. وأيونات الكلوريد تكسر
الغشاء عند نقاط ضعفه؛ فيكون الفلز العاري هناك مصعدًا صغيرًا أمام
مهبط مخمَّل هائل، فتكون كثافة تياره ضخمة. وداخل النقرة
تتحلمه أيونات الفلز وتحمّض المحلول، ويُجتذب الكلوريد إليها لموازنة
الشحنة، فلا يمكن للغشاء أن يتكوّن من جديد: فتتعمق النقرة.
\end{solution}

\begin{solution}{exo:b2:corrosion:11}
يُختزل الأكسجين على كل \qty{22}{cm^2}: \omterm{def:b2:current-potential-curves:ie-curve}{التيار المهبطي} $22 \times 5 =
\qty{110}{\micro A}$، يمده كله تأكسد \qty{2}{cm^2} من الفولاذ
قرب الوصلة، $j = \qty{55}{\micro A/cm^2}$. والفقد $5.5$ أضعاف فقد الحديد
عند \qty{10}{\micro A/cm^2}، أي نحو \qty{0.64}{mm} في السنة.
\end{solution}

\begin{solution}{exo:b2:corrosion:12}
يكون الحديد منيعًا تحت خط $\ce{Fe^2+}/\ce{Fe}$، $-0.41 + 0.0296\log 10^{-6} =
\qty{-0.59}{V}$: فإذا حُفظ تحت نحو \qty{-0.6}{V} لا يعود الفولاذ يذوب. وأدنى من ذلك
بكثير، يُختزل الماء على الفولاذ بسرعة متزايدة (فخط $\ce{H+}/\ce{H2}$ عند
\qty{-0.41}{V} عند pH 7 \omterm{def:b2:current-potential-curves:overpotential}{وفرط الجهد} على الفولاذ بضعة أعشار الفولط):
فيضيع التيار ويرفع الهيدروجين الطلاء.
\end{solution}

\begin{solution}{pb:b2:corrosion:1}
\textbf{1.} \ce{Fe -> Fe^2+ + 2e-}؛ \ce{O2 + 2H2O + 4e- -> 4OH-}.
\textbf{2.} اختزال الأكسجين، البطيء والبالغ استواء انتشاره،
يحد \omterm{def:b2:current-potential-curves:ie-curve}{التيار المهبطي}؛ ويلاقيه المنحنى المصعدي للحديد على ذلك الاستواء.
\textbf{3.} $j = \qty{0.10}{A/m^2}$: $0.10 \times 3.156 \times 10^7/(2 \times 96\,485)
\times 55.845 = \qty{913}{g}$ لكل متر مربع وفي السنة.
\textbf{4.} $913/7.87 \times 10^6 = \qty{1.16e-4}{m}$، أي \qty{0.12}{mm} في السنة.
\textbf{5.} $4/0.116 = 34$~سنة.
\textbf{6.} التآكل غير منتظم: فالشقوق، واختلافات التهوية
بين الترب، وعيوب الطلاء تركّز \omterm{def:b2:current-potential-curves:ie-curve}{التيار المصعدي} على
مساحات صغيرة، حيث تثقب النُّقَر الجدار قبل ذلك بكثير.
\textbf{7.} $E^\circ = \Delta_f G^\circ/(nF)$ لكل مزدوجة: \ce{Fe} $-0.41$، \ce{Zn} $-0.76$،
\ce{Mg} $-2.36$، \ce{Al} \qty{-1.68}{V}.
\textbf{8.} الزنك والمغنيسيوم والألومنيوم، وكلها تحت الحديد.
\textbf{9.} يتخمَّل الألومنيوم: فغشاء أكسيده يوقف ذوبانه ولا
يعطي عندئذ إلا تيارًا قليلًا، ما لم يُسبك بما يمنع الغشاء (وهو ما ينجح
في ماء البحر، لا في الترب).
\textbf{10.} يجب أن يعبر التيار مقاومة التربة: والجهد الدافع
بين المصعد والفولاذ المحمي نحو \qty{1.9}{V} للمغنيسيوم
مقابل \qty{0.35}{V} للزنك (من الجهود المعيارية).
\textbf{11.} $3.156 \times 10^7 \times 24.305/(2 \times 96\,485 \times 0.50) =
\qty{7.95}{kg}$ لكل أمبير وفي السنة.
\textbf{12.} $\pi \times 0.30 \times 10\,000 = \qty{9.42e3}{m^2}$.
\textbf{13.} \qty{94.2}{m^2}.
\textbf{14.} $0.020 \times 94.2 = \qty{1.88}{A}$.
\textbf{15.} $1.885 \times 20 \times 3.156 \times 10^7 = \qty{1.19e9}{C}$.
\textbf{16.} $1.19 \times 10^9 \times 24.305/(2 \times 96\,485 \times 0.50) =
\qty{3.00e5}{g}$، أي \qty{300}{kg} (أو $1.885 \times 20 \times 7.95$).
\textbf{17.} $300/15 = 20$ مصعدًا.
\textbf{18.} مقاومة التربة تحد المسافة التي يستطيع مصعد واحد أن يدفع إليها
تياره على طول الأنبوب: فالمصاعد الموزعة تحميه بانتظام.
\textbf{19.} يدفع مقوِّم التيار من مصعد خامل مدفون في فراش
مصعدي، عبر التربة، إلى الأنبوب الموصول بقطبه السالب.
\textbf{20.} $1.885 \times 5.0 + 1.5 = \qty{10.9}{V}$.
\textbf{21.} $10.9 \times 1.885 = \qty{20.6}{W}$؛ $20.6 \times 8766 = \qty{1.8e5}{W.h}$،
أي \qty{180}{kW.h} في السنة.
\textbf{22.} يُختزل الماء عند الأنبوب: فيتكوّن الهيدروجين، ويرفع الطلاء
وقد يقصف الفولاذ، ويضيع التيار.
\textbf{23.} $E < -0.41 + 0.0296\log(10^{-6}) = \qty{-0.59}{V}$.
\textbf{24.} تحتاج عشرون سنة من الحماية إلى $\boldsymbol{\approx \qty{3.0e2}{kg}}$
من مصاعد المغنيسيوم.
\end{solution}
